\section{Ranking System}

The ranking system is a core component of the platform, designed to motivate users, reflect individual learning progress, and encourage healthy competition among participants. It provides a transparent and consistent mechanism for evaluating user performance based on problem-solving activities.

\subsection*{Overall Ranking Mechanism}

The platform maintains a global leaderboard that ranks all registered users based on their accumulated points obtained from solving programming problems. This leaderboard reflects the relative standing of users across the entire system and allows learners to compare their progress with others. However, to reduce server workload and avoid unnecessary real-time synchronization, the ranking data is not pushed to clients in real time.

Instead, the frontend retrieves ranking data through the backend API when users access or refresh the leaderboard page. The backend calculates or returns the latest available ranking data based on accepted submissions and accumulated user scores. This request-based approach keeps the implementation simple and efficient while still providing users with sufficiently updated ranking information during normal platform usage.

Although the leaderboard is not updated through real-time events, users can view their current score and ranking information whenever the relevant page or profile data is loaded from the backend. This design provides a practical balance between system performance, implementation complexity, and user experience.

\subsection*{Point Allocation per Problem}

Points are awarded based on the difficulty level of each problem. The platform categorizes problems into three difficulty levels, each associated with a predefined base score:
\begin{itemize}
    \item \textbf{Easy}: 100 points
    \item \textbf{Medium}: 150 points
    \item \textbf{Hard}: 200 points
\end{itemize}

This scoring scheme encourages users to challenge themselves with more complex problems while still recognizing achievements at all difficulty levels.

\subsection*{Penalty for Failed Submissions}

To promote code quality and thoughtful problem-solving, the system applies a penalty for unsuccessful submissions. A submission is considered failed if it does not pass all predefined test cases. For each failed submission attempt, a penalty of 10 points is deducted from the base score of the corresponding problem.

If the total penalty incurred from failed submissions exceeds the base score of the problem, the final score awarded for that problem is capped at 0 points. Negative scores are not applied. This mechanism discourages excessive trial-and-error submissions while still allowing users to continue practicing without additional penalties once the minimum score is reached.

\subsection*{First Successful Submission Policy}

Points are awarded only for the first successful submission of a problem, defined as the first submission that passes all test cases. Any subsequent successful submissions for the same problem are treated strictly as practice attempts and do not contribute additional points to the user’s total score.

This policy ensures fairness in the ranking system and prevents users from inflating their scores by repeatedly submitting solutions to the same problem. It also emphasizes learning progression through solving new problems rather than optimizing repeated submissions for ranking purposes.
